\section{Conclusion}
\label{sec:conclusion}

We presented the first systematic characterization of how multi-class annotation errors interact with causal DAG topology. Our bias taxonomy across seven topologies reveals that the classical binary guarantee of bias toward the null breaks for $K \geq 3$: exposure and IV topologies can produce sign reversal (83.4\% and 52.6\% respectively), while five topologies maintain a structural guarantee of no sign flip. The ASV framework provides topology-aware sensitivity diagnostics, revealing a vulnerability spectrum from IV (maximally fragile) to M-bias (immune). Our correction guidance analysis reveals that correction method effectiveness is itself topology-dependent: a Frisch--Waugh--Lovell invariance result proves matrix-based corrections are ineffective for four topologies, while PPI++ achieves 93.8\% bias reduction at the cost of gold-standard labels. Binarization collapse analysis further demonstrates that simplifying multi-class annotations to binary introduces additional bias varying 147$\times$ across topologies.

\paragraph{Limitations.}
Our closed-form bias expressions are derived for $K{=}3$ and verified numerically at $K{=}5,7$; closed-form expressions for arbitrary $K$ remain open. The analysis assumes non-differential misclassification; differential error (where misclassification rates depend on the outcome) is more realistic for LLM annotators but substantially harder to characterize. Formal sharpness proofs for the monotonicity theorem cover three topologies; the remaining four are validated numerically (Appendix~\ref{app:monotonicity}). Sensitivity to DGP parameters is reported in Appendix~\ref{app:dgp_sensitivity}. The empirical validation covers three CSS tasks but uses analytical bias computation rather than end-to-end causal estimation on real outcome data.

\paragraph{Future work.}
Four directions are immediate. First, extending the framework to differential misclassification, which arises when LLM annotation quality varies with context. Second, incorporating the correlated error structure documented by \citet{kim_garg2025} into multi-LLM correction strategies. Third, extending the correction guidance framework to differential misclassification settings where PPI++ and DSL may interact differently with topology. Fourth, developing a practitioner toolkit that computes ASVs from estimated confusion matrices, routes to the appropriate correction method based on the researcher's causal design, and provides confidence intervals via PPI++ when gold-standard labels are available.
